% =============================================================================
% Abstract
% =============================================================================
\chapter*{Zusammenfassung}
\addcontentsline{toc}{chapter}{Zusammenfassung}

Die Röntgenabsorptionsspektroskopie (XAS), insbesondere die Analyse der
feinen Struktur oberhalb der Absorptionskante (EXAFS), liefert elementspezifische
Informationen über die lokale atomare Umgebung in kondensierten Materialien.
Bei Systemen mit mehreren chemischen Komponenten -- etwa ternären Legierungen
oder bimetallischen Nanopartikeln -- stößt die konventionelle Kurvenanpassung
jedoch an ihre Grenzen, da die Zahl der zu bestimmenden Strukturparameter
rasch die durch die Daten verfügbare Information übersteigt.

Die vorliegende Arbeit setzt die Software \evax{} ein, die auf einem
evolutionären Algorithmus in Kombination mit Reverse Monte Carlo (RMC)
basiert, um dieses Problem zu umgehen.
\evax{} erzeugt dreidimensionale Strukturmodelle und optimiert die atomaren
Positionen so, dass die simulierten EXAFS-Spektren aller beteiligten Kanten
simultan mit den experimentellen Daten übereinstimmen.

Im Mittelpunkt stehen zwei Proben:
eine \TODO{CoNiCu}-Mittelentropielegierung (MEA), deren drei K-Kanten
gleichzeitig ausgewertet werden, sowie \TODO{CoPt}-Nanopartikel, bei denen
die Verteilung der Elemente auf die Gitterplätze untersucht wird.
Für beide Systeme werden radiale Verteilungsfunktionen, Koordinationszahlen
und die chemische Nahordnung (Warren-Cowley-Parameter) bestimmt.
Darüber hinaus wird durch systematische Variation der Fit-Parameter und
den Vergleich unterschiedlicher Ausgangsstrukturen (FCC, BCC, HCP)
die Robustheit der Ergebnisse überprüft.

\vspace{0.5\baselineskip}

\textbf{Stichworte:}
EXAFS, Reverse Monte Carlo, evolutionärer Algorithmus, \evax{},
lokale Struktur, chemische Nahordnung, Mittelentropielegierung, Nanopartikel
